\documentclass{exam-zh}
\usepackage{siunitx}
\usepackage{draftwatermark} 

\SetWatermarkText{工试群版权所有，仅供学习，严禁商用}           % 设置水印内容
\SetWatermarkLightness{0.9}             % 设置水印透明度 0-1
\SetWatermarkScale{0.3}                   % 设置水印大小 0-1 


\examsetup{
  page/size=a4paper,
  paren/show-paren=true,
  paren/show-answer=true,
  fillin/show-answer=false,
  solution/show-solution=show-stay,
  solution/label-indentation=false,
  page/foot-content={试题第;页（共;页）}
}

\everymath{\displaystyle}

\title{上海交通大学2024年平台预分专业能力测试考试}

\subject{数学和逻辑}



\begin{document}

% \information{
%   姓名\underline{\hspace{6em}},
%   座位号\underline{\hspace{15em}}
% }
% \warning{（在此卷上答题无效）}

% \secret

\maketitle

本试卷共 60 题。全卷满分 100 分。考试用时 120 分钟。


\begin{notice}
  \item 上海交通大学工科试验班24群（QQ 760990141）版权所有，禁止商用
  \item 制作人员：方天{\CJKfontspec{华文宋体}宬}（策划/整理/排版）、铁镐镐（整理）、w1049、谭宇、方解石、杜大佬、24-方cs-pluto、张光迪、王浩宇、袁银豪、李亚、项俊豪、BB不是BBQ、日月星光、王鸿恺、梁铠瑜、KiloxGo、许驰声、粥粥、蟹老板、youxi、刘法样、Pollux、郭颖翰、HugoYu、T、Mario、樱吹雪、空诗、Qy、do-ne、一条三文鱼、Prop.、yuzusofter、朝比奈真冬、贺耀、陈天
  \item 本卷是回忆版，疏漏之处请见谅
  \item 温馨提示：大多数考生在规定时间内只完成了50题左右
\end{notice}



\section{%
  数学部分：共 30 小题，1 到 10 题每小题 1 分，11 到 30 题每小题 2 分。
}

%1
\begin{question}
  二项式$\left(2x^3-\frac1{\sqrt{x}}\right)^7$展开式中的常数项为 \paren[]

  \begin{choices}
    \item -14
    \item -7
    \item 7
    \item 14
  \end{choices}
\end{question}

%2
\begin{question}
  设$a$和$b$是方程$x^2+3x+7=0$的两个根，则$a^3+b^3=$ \paren

  \begin{choices}
    \item -36
    \item -19
    \item 19
    \item 36
  \end{choices}
\end{question}

%3
\begin{question}
  已知等比数列$\{a_n\}$的前$n$项和为$S_n$，则下列一定成立的是 \paren

  \begin{choices}
    \item $a_3>0$，则$a_{2024}>0$
    \item $a_4>0$，则$a_{2023}>0$
    \item $a_3>0$，则$S_{2023}>0$
    \item $a_4>0$，则$S_{2024}>0$
  \end{choices}
\end{question}

% 4.
\begin{question}
  在$\triangle ABC$中，$\angle A=\frac\pi3$，$\angle B=\frac\pi4$，$BC=2\sqrt2$，则$AC=$ \paren

  \begin{choices}
    \item $4\sqrt3$
    \item $2\sqrt3$
    \item $\sqrt3$
    \item $\frac{\sqrt3}2$
  \end{choices}
\end{question}

% 5.
\begin{question}
  在极坐标系中$\theta=0$，$\theta=\frac\pi3$，$r\sin\theta+r\cos\theta=1$围成图形的面积是 \paren

  \begin{choices}
    \item $\frac{3+\sqrt3}4$
    \item $\frac{3-\sqrt3}4$
    \item $\frac{\sqrt3+1}4$
    \item $\frac{\sqrt3-1}4$
  \end{choices}
\end{question}

% 6.
\begin{question}
  $\alpha,\beta$是方程$\lg^2 x-\lg x^2-3=0$的两根，则$\log_\alpha\beta+\log_\beta\alpha=$ \paren

  \begin{choices}
    \item $-\frac{10}3$
    \item $-2$
    \item $0$
    \item $1$
  \end{choices}
\end{question}

% 7.
\begin{question}
  复数$z\in\mathbb C$满足$|z-1|=1$，则$|z+2+4i|$的最小值为 \paren

  \begin{choices}
    \item $\sqrt{17}-1$
    \item 4
    \item $\sqrt{17}+1$
    \item 6
  \end{choices}
\end{question}

% 8.
\begin{question}
  函数$f(x)=\ln x$， 对$x_1,x_2>0$且$x_1\neq x_2$，以下说法正确的个数是 \paren

  \begin{enumerate}
    \item $f(x_1+x_2)=f(x_1)+f(x_2)$
    \item $f(x_1x_2)=f(x_1)+f(x_2)$
    \item $\frac{f(x_1)-f(x_2)}{x_1-x_2}>0$
    \item $f\left(\frac{x_1+x_2}2\right)>\frac{f(x_1)+f(x_2)}2$
  \end{enumerate}

  \begin{choices}
    \item 1
    \item 2
    \item 3
    \item 4
  \end{choices}
\end{question}

% 9.
\begin{question}
  已知圆$C$的方程为$x^2+(y-2)^2=a$，则“$a>2$”是“函数$y=|x|$的图像与圆$C$有四个公共点”的 \paren

  \begin{choices}
    \item 充分不必要条件
    \item 必要不充分条件
    \item 充要条件
    \item 既不充分也不必要条件
  \end{choices}
\end{question}

% 10.
\begin{question}
  已知数据$x_1,x_2,\cdots,x_{10}$满足：$x_i-x_{i-1}=1\,(2\le i\le10)$。，去掉$x_1$和$x_{10}$组成一组新数据，则新数据与原数据相比，下列说法错误的是 \paren

  \begin{choices}
    \item 中位数不变
    \item 平均数不变
    \item 方差变小
    \item 若$x_1=1$，则数据$x_1,\cdots,x_{10}$的75\%分位数是7.5
  \end{choices}
\end{question}

% 11.
\begin{question}
  已知$\alpha\in\left(0,\frac\pi2\right)$，$\beta\in\left(0,\frac\pi2\right)$，且$\sin\left(2\alpha+\beta\right)=\frac32\sin\beta$，则$\cos\beta$的最小值为 \paren

  \begin{choices}
    \item $\frac{\sqrt 5}3$
    \item $\frac{\sqrt 5}5$
    \item $\frac12$
    \item $\frac23$
  \end{choices}
\end{question}

% 12.
\begin{question}
  定义在$\mathbb R$上的函数$f(x)=\frac{4^x}{4^x+2}$，则$\sum_{n=1}^{2023}f\left(\frac{n}{2024}\right)=$ \paren

  \begin{choices}
    \item $\frac{2023}2$
    \item 1012
    \item $\frac{2025}2$
    \item 1013
  \end{choices}
\end{question}

% 13.
\begin{question}
  数列$\{x_n\}$满足$x_{n+1}=\frac{2(1+x_n)}{2+x_n},\;n=1,2,3$，且$x_1\neq1$，$x_1\neq\sqrt2$，则 \paren

  \begin{choices}
    \item $\{x_n\}$单调增加
    \item $\{x_n\}$可能单调增加，也可能单调减小
    \item $\{x_n\}$单调减小
    \item $\{x_n\}$非单调增，也非单调减
  \end{choices}
\end{question}

% 14.
\begin{question}
  $\mathbb R^3$空间中，$O$发出四条射线，两两间夹角相同，则该角度为 \paren

  \begin{choices}
    \item $\frac\pi2$
    \item $\frac23\pi$
    \item $2\arctan\sqrt2$
    \item $2\arctan2$
  \end{choices}
\end{question}

% 15.
\begin{question}
  定义在$n\in \mathbb N_+$的函数$f(n)=3^{n+1}(2n+9)+9$满足$f(n)$总是$m$的倍数，则$m$最大为 \paren

  \begin{choices}
    \item 9
    \item 18
    \item 27
    \item 36
  \end{choices}
\end{question}

% 16.
\begin{question}
  在立方体中随机选三条棱，两两不共面的概率为 \paren

  \begin{choices}
    \item $\frac1{55}$
    \item $\frac2{55}$
    \item $\frac3{55}$
    \item $\frac1{11}$
  \end{choices}
\end{question}


% 17.
\begin{question}
  在平面直角坐标系$xOy$中，$K=\left\{(x,y)|x,y=-1,0,1\right\}$，在$K$中随机取2个点，其距离是$\sqrt5$的概率是 \paren

  \begin{choices}
    \item $\frac19$
    \item $\frac29$
    \item $\frac13$
    \item $\frac49$
  \end{choices}
\end{question}

% 18.
\begin{question}
  所有形如“$1010\dots101$”，首位和末尾都是$1$，且相邻位数$0$和$1$交替出现的自然数中，素数有 \paren

  \begin{choices}
    \item 1个
    \item 2个
    \item 3个
    \item 无穷个
  \end{choices}
\end{question}

% 19.
\begin{question}
  已知圆锥的顶点为$S$，$O$为底面圆心，母线$SA$与$SB$垂直，$\triangle SAB$的面积为$2$，$PA$与底面夹角为$\frac\pi6$，下列正确的是 \paren

  \begin{choices}
    \item 圆锥高为$\sqrt3$
    \item 圆锥侧面展开扇形圆心角为$\frac23\pi$
    \item 面$SAB$与底面夹角为$\frac\pi4$
    \item 圆锥侧面积为$3\pi$
  \end{choices}
\end{question}

% 20.
\begin{question}
  椭圆两焦点$F_1F_2$位于$x$轴，上顶点为$A$，$AF_2$交椭圆于另一点$B$，且$|AF_2|=2|F_2B|$，则椭圆的离心率等于 \paren

  \begin{choices}
    \item $\frac{\sqrt 3}2$
    \item $\frac{\sqrt 3}3$
    \item $\frac{\sqrt 2}2$
    \item $\frac{\sqrt 2}3$
  \end{choices}
\end{question}

% 21.
\begin{question}
  $O$为坐标原点，双曲线$C$的方程为$x^2-\frac{y^2}3=1$，直线$y=x+2$与$C$的渐近线交于$AB$两点，求$\triangle OAB$内切圆半径 \paren

  \begin{choices}
    \item $\sqrt 2-1$
    \item $2-\sqrt3$
    \item $2-\sqrt2$
    \item $\sqrt3-1$
  \end{choices}
\end{question}

% 22.
\begin{question}
  甲抛2次硬币，乙抛3次，事件$M$为“甲第一次抛硬币结果为正面”，事件$N$为“甲抛硬币两次结果相同”，事件$S$为“甲正面的次数少于乙”，则 \paren

  \begin{choices}
    \item $M\subseteq N$
    \item $P(M\cup N)=P(M)+P(N)$
    \item $P(S)<P(N)$
    \item $P(S)=P(M)$
  \end{choices}
\end{question}

% 23.
\begin{question}
  已知4个数的平均数是5，方差是4，加入6和10后，方差变为 \paren

  \begin{choices}
    \item $\frac{13}3$
    \item $\frac53$
    \item 6
    \item 8
  \end{choices}
\end{question}

% 24.
\begin{question}
  在$\triangle ABC$中，内角$ABC$对应的边分别为$abc$，三角形面积为$S$，若$\sin C=\frac{2S}{a^2+b^2-c^2}$，且$\triangle ABC$外接圆半径为$\frac23\sqrt{3}$，求$a+b+c$的取值范围为 \paren

  \begin{choices}
    \item $(4,6]$
    \item $(4,4\sqrt3]$
    \item $(6,8]$
    \item $(6,6\sqrt 3]$
  \end{choices}
\end{question}

% 25.
\begin{question}
  $P$为棱长1立方体$ABCD-A'B'C'D'$表面上一点，且$\vec{AP}\cdot\vec{BD_1}=1$，求$|AP|$的取值范围为 \paren

  \begin{choices}
    \item $[1,\sqrt3]$
    \item $[\frac{\sqrt 2}2,\sqrt3]$
    \item $[1, \sqrt2]$
    \item $[\frac12,\sqrt2]$
  \end{choices}
\end{question}

% 26.
\begin{question}
  函数$f(x)=\cos\pi x-|\log_2|x-1||$的所有零点之和为 \paren

  \begin{choices}
    \item 0
    \item 2
    \item 4
    \item 6
    \item 8
  \end{choices}
\end{question}

% 27.
\begin{question}
  函数$f(x)=\begin{cases}\frac ex\ln x,&x>0\\-2024x,&x\le0\end{cases}$，$g(x)=f^2(x)-(2m-1)x+2$有4个零点，求$m$的取值范围 \paren

  \begin{choices}
    \item $m>2$
    \item $m\ge2$
    \item $m>\frac{\sqrt2+1}2$或$m<\frac{\sqrt2-1}2$
    \item $m>\frac{\sqrt2+1}2$
  \end{choices}
\end{question}

% 28.
\begin{question}
  $(2x-3)^9=a_0+a_1(x-1)+a_2(x-1)^2+\cdots+a_9(x-1)^9$，则$a_0+2a_1+3a_2+\cdots+10a_9=$ \paren

  \begin{choices}
    \item 9
    \item 10
    \item 19
    \item 29
  \end{choices}
\end{question}

% 29.
\begin{question}
  函数$f(x)=\frac1x-x+a\ln x$，判断以下两个命题的正确性 \paren

  \begin{enumerate}
    \item $a\le2$时，$f(x)$在$(0,+\infty)$单调递减
    \item $f(x)$有两个极值点$x_1,x_2$时，$\frac{f(x_1)-f(x_2)}{x_1-x_2}<a-2$
  \end{enumerate}
  \begin{choices}
    \item ① 正确，② 错误
    \item ① 错误，② 正确
    \item ① 正确，② 正确
    \item ① 错误，② 错误
  \end{choices}
\end{question}

% 30.
\begin{question}
  $abc=1$，$a,b,c>0$，且$a,b,c\neq1$，求$\frac{a^3}{b^2(1-b)(a-1)}+\frac{b^3}{c^2(1-c)(b-1)}+\frac{c^3}{a^2(1-a)(c-1)}$的最小值为 \paren

  \begin{choices}
    \item 10
    \item 18
    \item 20
    \item 28
  \end{choices}
\end{question}

\section{%
  逻辑部分：共 30 小题，31 到 40 题每小题 1 分，41 到 60 题每小题 2 分。缺失部分跳过，部分题目题号遗忘，从 61 开始编号。
}

% 31.
\begin{question}
  博物馆规定，个人票每张10元，团体票每张60元（可供10人参观），无其他票价优惠政策。共有58人逛A公园，则最少应付多少元? \paren

  \begin{choices}
    \item 350
    \item 360
    \item 380
    \item 390
  \end{choices}
\end{question}

% 32.
\begin{question}
  某商品按定价出售，每个可获得60元的利润。按定价打八折出售10个所获得的利润，与按定价每个减价30元出售15个所获得的利润相同。该商品的定价为多少元？ \paren
  
  \begin{choices}
    \item 75
    \item 80
    \item 85
    \item 90
  \end{choices}
\end{question}

% 33.
\begin{question}
  小张去年底获得一笔总额不超过5万的奖金，她将其中的60\%用来储蓄，剩下的用来购买理财产品，一年后这笔奖金增值了5\%。已知储蓄的奖金增值了3.3\%，问购买理财产品的奖金增值了多少？ \paren
  
  \begin{choices}
    \item 5.35\%
    \item 6.45\%
    \item 7.55\%
    \item 8.65\%
  \end{choices}
\end{question}

% 34.
\begin{question}
  学生说：“有了努力不一定能成功考上研究生。”老师鼓励说：“但是没有努力一定会失败。”以下哪项与老师的意思最为接近? \paren
  
  \begin{choices}
    \item 不成功即失败，不失败即成功。
    \item 如果努力，则一定会成功。
    \item 只有努力，才可能不失败。
    \item 除非努力，否则不可能失败。
    \item 只有成功了，才可能更努力。
  \end{choices}
\end{question}

% 35.
\begin{question}
  一个人如果没有崇高的信仰，就不可能守住道德的底线，而一个人只有不断加强理论学习，才能始终保持崇高的信仰。根据以上信息，可以得出以下哪项？ \paren

  \begin{choices}
    \item 一个人没能守住道德的底线，是因为他首先丧失了崇高的信仰
    \item 一个人只要有崇高的信仰，就能守住道德的底线
    \item 一个人只有不断加强理论学习，才能守住道德的底线
    \item 一个人如果不能守住道德的底线，就不可能保持崇高的信仰
    \item 一个人只要不断加强理论学习，就能守住道德的底线
  \end{choices}
\end{question}

% 36.
\begin{question}
  如果他勇于承担责任，那么他就一定会直面媒体，而不是选择逃避；如果他没有责任，那么他就一定会聘请律师，捍卫自己的尊严。可是事实上，他不仅没有聘请律师，现在逃得连人影都不见了。根据以上陈述，可以得出以下哪项结论？ \paren

  \begin{choices}
    \item 即使他没有责任，也不应该选择逃避。
    \item 虽然选择了逃避，但是他可能没有责任。
    \item 如果他有责任，那么他应该勇于承担责任。
    \item 如果他不敢承担责任，那么说明他责任很大。
    \item 他不仅有责任，而且他没有勇气承担责任。
  \end{choices}
\end{question}

\questionsetup{index=42}
% 42.
\begin{question}
  甲、乙两个单位周末分别安排60\%和75\%的职工下沉社区帮助困难群众，其中甲单位派出的职工比乙单位少3人。后两单位又在剩下的职工中，分别抽调40\%和75\%的职工，共计24人参加周末的业务培训。问甲单位职工人数比乙单位： \paren

  \begin{choices}
    \item 少 3 人
    \item 少 11 人
    \item 多 3 人
    \item 多 11 人
  \end{choices}
\end{question}

% 43.
\begin{question}
  甲乙两地铁路线长1880千米，从甲地到乙地开出一辆动车，每小时行驶160千米，3小时后，从乙地到甲地开出一辆高铁，经4小时后与动车相遇，则高铁每小时行驶： \paren

  \begin{choices}
    \item 180千米
    \item 210千米
    \item 200千米
    \item 190千米
  \end{choices}
\end{question}

\questionsetup{index=46}
% 46.
\begin{question}
  甲、乙两个箱子中分别装有不同数量的某种商品，总数不到100件。如果从甲或乙箱子中随机拿出1件这种商品，拿到次品的概率分别为5\%和10\%。如果将两箱内的商品混合后再随机拿出1件，则拿到次品的概率为6\%，从乙箱中随机拿出3件这种商品，均不为次品的概率$x$在以下哪个范围内？ \paren

  \begin{choices}
    \item $x < 72\%$
    \item $72\% \le x < 81\%$
    \item $81\% \le x < 90\%$
    \item $x \ge 90\%$
  \end{choices}
\end{question}

\questionsetup{index=48}
% 48.
\begin{question}
  1、2、3、4、5、6、7、8、9这九个数字各用一次，组成三个能被9整除的三位数，这三个数的和最大是： \paren

  \begin{choices}
    \item 2007
    \item 2394
    \item 2448
    \item 2556
    \item 2635
  \end{choices}
\end{question}

% 49.
\begin{question}
  某单位举办两轮羽毛球单打表演赛，共有甲、乙、丙、丁、戊、己6位选手参加。每轮表演赛都按以下组合进行了5场比赛：甲对乙、甲对丁、丙对戊、丙对丁、戊对己。已知：
  \begin{enumerate}
    \item 每场比赛均决出胜负；
    \item 每轮比赛中，各参赛选手均至多输一场；
    \item 每轮比赛决出的冠军在该轮比赛中未有败绩，甲在第一轮比赛中获冠军；
    \item 只有一组选手在第二轮比赛中的胜负结果与第一轮相同，其余任一组选手的两轮比赛结果均不同。
  \end{enumerate}
  根据上述信息，可以得出第二轮表演赛的冠军是？ \paren

  \begin{choices}
    \item 乙
    \item 丙
    \item 丁
    \item 戊
    \item 己
  \end{choices}
\end{question}

% 50.
\begin{question}
  张霞、李丽、陈露、邓强和主硕一起坐火车去旅游。他们正好坐在同一车厢相对两排的五个座位上，每人各坐一个位置。第一排的座位按顺序分别记作1号和2号，第2排的座位按顺序记为3、4、5号。座位1和座位3直接相对，座位2和座位4直接相对。座位5不和上述任何座位直接相对。李丽坐在4号位置，陈露所坐的位置不与李丽相邻，也不与邓强相邻（相邻是指同一排上紧挨着），张霞不坐在陈露直接相对的位置上。根据以上信息，张霞所坐位置有多少种可能的选择？ \paren

  \begin{choices}
    \item 5种
    \item 4种
    \item 3种
    \item 2种
    \item 1种
  \end{choices}
\end{question}

% 51.
\begin{question}
  某次讨论会共有18名参会者。已知：\begin{enumerate}
    \item 至少有5名青年教师是女性；
    \item 至少有6名女教师已过中年；
    \item 至少有7名女青年是教师。
  \end{enumerate}
  根据上述信息，关于参会人员可以得出以下哪项？ \paren

  \begin{choices}
    \item 有些青年教师不是女性
    \item 有些女青年不是教师
    \item 青年教师至少有11名
    \item 女青年至多有11名
    \item 女教师至少有13名
  \end{choices}
\end{question}

% 52.
\begin{question}
  某次讨论会共有18名参会者。已知：\begin{enumerate}
    \item 至少有5名青年教师是女性；
    \item 至少有6名女教师已过中年；
    \item 至少有7名女青年是教师。
  \end{enumerate}
  如果上述三句话两真一假，那么关于参会人员可以得出以下哪项？ \paren

  \begin{choices}
    \item 青年教师至少有5名
    \item 男教师至多有10名
    \item 女青年都是教师
    \item 女青年至少有7名
    \item 青年教师都是女性
  \end{choices}
\end{question}

\questionsetup{index=60}

% 60.
\begin{question}
  某大学的哲学学院和管理学院今年招聘新教师，招聘结束后受到了女权主义代表的批评，因为他们在12名女性应聘者中录用了6名，但在12名男性应聘者中却录用了7名。该大学对此解释说，今年招聘新教师的两个学院中，女性应聘者的录用率都高于男性应聘者的录用率。具体情况是：哲学学院在8名女性应聘者中录用了3名，而在3名男性应聘者中录用了1名；管理学院在4名女性应聘者中录用了3名，而在9名男性应聘者中录用了6名。以下哪项最有助于解释女权主义代表和大学之间的分歧？ \paren

  \begin{choices}
    \item 整体并不是局部的简单相加
    \item 有些数字规则不能解释社会现象
    \item 人们往往从整体角度考虑问题，不管局部
    \item 现代社会提倡男女平等，但实际执行中还有一定难度
    \item 各个局部都具有的性质在整体上未必具有
  \end{choices}
\end{question}

% 题号未知
\begin{question}
  一杯质量分数为15\%的200g盐水，需要加多少克盐，才能使之变为质量分数为20\%的盐水？ \paren

  \begin{choices}
    \item 12.5
    \item 10
    \item 5.5
    \item 5
  \end{choices}
\end{question}

% 题号未知
\begin{question}
  某大学从候选人甲、乙、丙、丁、戊、己、庚7人中选出3人作为本年度优秀教师。已知：\begin{enumerate}
    \item 甲、丙、丁、戊、己中至多有2人入选；
    \item 若戊、己都没有入选，则丁、庚也都没有入选；
    \item 若乙、庚中至少有1人没入选，则甲、丙都入选。
  \end{enumerate}
  根据上述信息，可以得出以下哪项？ \paren

  \begin{choices}
    \item 甲入选
    \item 乙入选
    \item 丙入选
    \item 戊入选
    \item 庚入选
  \end{choices}
\end{question}

% 题号未知
\begin{question}
  陈先生在鼓励他的孩子时说道：“不要害怕暂时的困难与挫折，不经历风雨怎么见彩虹？”他的孩子不服气的说：“您说的不对，我经历了那么多风雨，怎么就没见到彩虹呢？”陈先生孩子的回答最适宜用来反驳以下哪项？ \paren

  \begin{choices}
    \item 只要经历了风雨，就可以见到彩虹
    \item 如果想见到彩虹，就必须经历风雨
    \item 只有经历风雨，才能见到彩虹
    \item 即使经历了风雨，也可能见不到彩虹
  \end{choices}
\end{question}

% 题号未知
\begin{question}
  公司经理：我们招聘人才时最看重的是综合素质和能力，而不是分数。人才招聘中，高分低能者并不鲜见，我们显然不希望招到这样的“人才”，从你的成绩单可以看出，你的学业分数很高，因此我们有点怀疑你的能力和综合素质。以下哪项和经理得出结论的方式最为类似？ \paren

  \begin{choices}
    \item 公司管理者并非都是聪明人，陈然不是公司管理者，所以陈然可能是聪明人
    \item 猫都爱吃鱼，没有猫患近视，所以吃鱼可以预防近视
    \item 人的一生中健康开心最重要，名利都是浮云，张立名利双收，所以可能张立并不开心
    \item 有些歌手是演员，所有的演员都很富有，所以有些歌手可能不富有
    \item 闪光的物体并非都是金子，考古队挖到了闪闪发光的物体，所以考古队挖到的可能不是金子
  \end{choices}
\end{question}

% 题号未知
\begin{question}
  3, 8, 23, 68, ( ), 608 \paren

  \begin{choices}
    \item 183
    \item 188
    \item 203
    \item 208
  \end{choices}
\end{question}

% 题号未知
\begin{question}
  比较文字学者张教授认为，在不同的民族语言中，字形与字义的关系有不同的表现。他提出，汉字是象形文字，其中大部分的是形声字，这些字的字形与字义相互关联；而英语是拼音文字，其字形与字义往往关联度不大，需要某种抽象的理解。以下哪项如果为真，最不符合张教授的观点？ \paren

  \begin{choices}
    \item 汉语中的“日”“月”是象形字，从字形可以看出其所指的对象；而英语中的 sun与moon 则感觉不到这种形义结合
    \item 汉语中的“日”与“木”结合，可以组成“束”“果”“杳”等不同的字，并可以猜测其语义。而英语中则不存在与此类似的 sun与 wood的结合
    \item 英语中，也有与汉语类似的象形文字，如 eye是人的眼睛的象形，两个e代表眼睛，y 代表中间的鼻子；bed是床的象形，b和d 代表床的两端
    \item 英语中的 sunlight与汉语中的“阳光”相对应，而英语的sun与light 和汉语中的“阳”与“光”相对应
    \item 汉语的“星期三”与英语中的 Wednesday和德语中的 MIttwoch意思相同
  \end{choices}
\end{question}

% 题号未知
\begin{question}
  某大学顾老师在回答有关招生问题时强调：“我们学校招收一部分免费师范生，也招收一部分一般师范生。一般师范生不同于免费师范生，没有免费师范生毕业时可以留在大城市工作，而一般师范生毕业时都可以选择留在大城市工作，任何非免费师范生毕业时都需要自谋职业，没有免费师范生毕业时需要自谋职业。”根据顾老师的陈述，可以得出以下哪项？ \paren

  \begin{choices}
    \item 该校需要自谋职业的大学生都可以选择留在大城市工作
    \item 该校可以选择留在大城市工作的唯一一类毕业生是一般师范生
    \item 不是一般师范生的该校大学生都是免费师范生
    \item 该校所有一般师范生都需要自谋职业
    \item 该校需要自谋职业的大学生都是一般师范生
  \end{choices}
\end{question}

% 题号未知
\begin{question}
  在某次考试中，有3个关于北京旅游景点的问题，要求考生每题选择某个景点的名字作为唯一答案。其中6位考生关于上述3个问题的答案依次如下：第一位考生：天坛、天坛、天安门；第二位考生：天安门、天安门、天坛；第三位考生：故宫、故宫、天坛；第四位考生：天坛、天安门、故宫；第五位考生：天安门、故宫、天安门；第六位考生：故宫、天安门、故宫。考试结果表明，每位考生都至少答对其中1道题。根据以上陈述，可知这三个问题的正确答案依次是：
  \paren

  \begin{choices}
    \item 天坛、故宫、天坛
    \item 故宫、天安门、天安门
    \item 天安门、故宫、天坛
    \item 天坛、天坛、故宫
    \item 故宫、故宫、天坛
  \end{choices}
\end{question}

% 题号未知
\begin{question}
  张老师说：这次摸底考试，我们班的学生全都通过了，所以，没有通过的都不是我们班的学生。以下哪项和以上推理最为相似？
  \paren

  \begin{choices}
    \item 所有摸底考试通过的学生都好好复习了，所以好好复习的学生都通过了
    \item 所有摸底考试没有通过的学生都没有好好复习，所以没有好好复习的学生都没有通过
    \item 所有参加摸底考试的学生都经过了认真准备，所以没有参加摸底考试的学生都没有认真准备
    \item 英雄都是经得起考验的，所以经不起考验的就不是英雄
    \item 有的学生虽然没有好好复习，但是也通过了摸底考试
  \end{choices}
\end{question}

% 题号未知
\begin{question}
  国际足联一直坚称，世界杯冠军队所获得的“大力神”杯是实心的纯金奖杯。某教授经过精密测量和计算认为，世界杯冠军奖杯——实心的“大力神”杯不可能是纯金制成的，否则球员根本不可能将它举过头顶并随意挥舞。以下哪项与这位教授的意思最为接近？
  \paren

  \begin{choices}
    \item 只有球员能够将“大力神”杯举过头顶并自由挥舞，它才是由纯金制成，并且不是实心的
    \item 若球员能够将“大力神”杯举过头顶并自由挥舞，则它很可能是空心的纯金杯
    \item 若“大力神”杯是实心的纯金杯，则球员不可能把它举过头顶并随意挥舞
    \item 只有“大力神”杯是实心的，它才可能是纯金的
    \item 若“大力神”杯是由纯金制成，则它肯定是空心的
  \end{choices}
\end{question}

% 题号未知
\begin{question}
  某电商网站推出免息分期购物活动，购买某件商品的消费者第一个月只用支付总金额的一半加10元，第二个月支付剩余金额的一半加20元，第三个月支付剩余金额的一半加30元，第四个月付清剩余未支付的10元。问这件商品的价格为多少元？
  \paren

  \begin{choices}
    \item 400
    \item 410
    \item 420
    \item 460
  \end{choices}
\end{question}

\end{document}